\section*{Situation Engineering framework}

Prompt engineering specifies a task; context and retrieval engineering select
information; memory engineering retains observations; tools expose
capabilities; and harness engineering controls execution. Situation engineering
addresses a narrower semantic question: which available claims apply to this
query, observer and world now? It is therefore a proposed interface among
existing components, not a replacement name for them and not an empirically
validated universal layer.

For query $q$, evidence $E$, retained memory $M_t$, tool observations $O_t$ and
harness trace $H_t$, situation construction produces candidate states
$\mathcal{S}_t$ with unresolved variables $U_t$ and provenance $P_t$:
\begin{equation}
(q,E,M_t,O_t,H_t)\rightarrow(\mathcal{S}_t,U_t,P_t).
\end{equation}
Each evidence claim is assigned a query-relative applicability value
\begin{equation}
\mathrm{Apply}(e,q,S_t)\in
\{\mathrm{valid},\mathrm{invalid},\mathrm{unresolved}\}.
\end{equation}
Preserving \emph{unresolved} is essential. A system should retrieve, clarify or
abstain when answer-critical state cannot be established rather than silently
converting uncertainty into a plausible binary judgement.

\begin{figure}[t]
\centering
\resizebox{\linewidth}{!}{\begin{tikzpicture}[
every node/.style={font=\small}, box/.style={draw,rounded corners,align=center,
minimum height=8mm}, arr/.style={-{Latex},thick}]
\node[anchor=west,font=\bfseries] at (0,4.5) {(a) Failure};
\node[box,text width=2.8cm] (appointed) at (3.0,4.5) {appointed at $t_1$};
\node[box,text width=2.8cm] (resigned) at (7.0,4.5) {resigned at $t_2$};
\node[box,text width=3.4cm,fill=red!8] (question) at (11.5,4.5) {Who is current at $t_3$?};
\draw[arr] (appointed)--(resigned); \draw[arr] (resigned)--(question);

\node[anchor=west,font=\bfseries] at (0,3.0) {(b) Retrieval only};
\node[box,text width=3.0cm] (claims) at (3.0,3.0) {relevant claims};
\node[box,text width=3.0cm] (rank) at (7.0,3.0) {rank / concatenate};
\node[box,text width=3.4cm,fill=red!8] (answer) at (11.5,3.0) {supported answer};
\draw[arr] (claims)--(rank); \draw[arr] (rank)--(answer);

\node[anchor=west,font=\bfseries] at (0,1.5) {(c) Situation engineering};
\node[box,text width=3.0cm] (input) at (3.0,1.5) {query, evidence,\\memory, tools};
\node[box,text width=3.0cm] (state) at (7.0,1.5) {sense + update\\provenance-linked state};
\node[box,text width=3.4cm,fill=green!10] (policy) at (11.5,1.5) {validate applicability\\answer / retrieve / clarify / abstain};
\draw[arr] (input)--(state); \draw[arr] (state)--(policy);

\node[anchor=west,font=\bfseries] at (0,0) {(d) Diagnostic axes};
\node[box,text width=11.0cm] at (7.2,0) {time \quad scope \quad modality \quad source conflict \quad observer \quad world \quad unresolved premise};
\end{tikzpicture}}
\caption{From situation blindness to an applicability-aware response.
\textbf{a}, A supported appointment becomes stale after resignation.
\textbf{b}, retrieval and relevance alone do not require resolution of that
transition. \textbf{c}, situation engineering constructs a provenance-linked
state, validates claim applicability and propagates unresolved variables into
response policy. \textbf{d}, SituationCatch-Bench isolates seven corresponding
diagnostic axes.}
\label{fig:engineering-planes}
\end{figure}

\subsection*{Operations and component boundaries}

Six operations define the interface. \emph{Sensing} extracts actors, events,
times, modality, source status, scope, observer knowledge and world identity.
\emph{Assembly} links observations into candidate states. \emph{Updating}
applies supersession, cancellation and belief transitions. \emph{Validation}
checks applicability, conflicts, provenance and missing variables.
\emph{Policy} selects an answer, retrieval, clarification, abstention or action.
\emph{Observability} records the active state and its evidence.

These responsibilities do not absorb neighbouring disciplines. A retriever
still owns ranking and freshness; a memory system owns retention; a tool layer
owns invocation semantics; a harness owns permissions and recovery; and serving
systems own resource efficiency. The proposed situation layer consumes their
outputs and returns applicability decisions and unresolved-state signals.
Likewise, longer context is not sufficient: it may contain an appointment and
a later resignation, a proposal and a final decision, or real and hypothetical
claims simultaneously. Applicability requires resolving their relation rather
than merely retaining both.

\subsection*{SCQA as a bounded control implementation}

Situation-Catching Question Answering (SCQA) instantiates the interface for
seven controlled variables. Its deterministic sensor or supplied gold state
feeds auditable update rules and an answer/clarify/abstain policy. SCQA has no
learned parameters. It is designed to test whether removing a required variable
produces the matching failure signature, not to compete with contemporary
language models.

This scope matters for interpretation. The generated benchmark and solver share
an ontology, making the oracle-state experiment analogous to a protocol
conformance suite. The predicted-state condition introduces sensing error but
still uses generated English. Natural-language generalization, comparison with
strong prompting and retrieval systems, multilingual transfer, human
annotation and deployment efficiency remain untested hypotheses. The released
multi-provider and annotation programs make these studies executable but do not
turn an unexecuted protocol into evidence.

\subsection*{Testable predictions}

The framework yields falsifiable predictions for future matched experiments.
Holding model, evidence and computation budget fixed, an explicit situation
state should improve the categories whose decisive variables are otherwise
implicit. Gains should diminish when a baseline already represents the same
state. Predicted-state errors should decompose into sensing, update, policy and
generation failures. Corrupting one state slot should selectively damage its
matching category. Finally, the benefit must be evaluated jointly with token
use, latency, retrieval calls, monetary cost and unnecessary clarification.
Failure of these predictions would argue against treating situation engineering
as a distinct practical interface.
